\chapter{Résultats et analyse}
\label{chap:resultats}

Ce chapitre présente les résultats obtenus en appliquant le modèle décrit au chapitre~\ref{chap:methodo} au profil opérationnel de Sabodala-Massawa. Sauf indication contraire, les montants sont exprimés en millions de dollars US constants (« réels »), cumulés sur la durée de vie du projet (2020-2036), conformément à la présentation retenue par l'outil NRGI dans sa feuille de résultats.

\section{Flux de trésorerie du projet et partage entre l'État et l'investisseur}

Le flux de trésorerie du projet avant toute imposition --- c'est-à-dire la différence entre le chiffre d'affaires et la somme des dépenses d'investissement, des coûts d'exploitation et des coûts de fermeture --- atteint 3\,595,4~millions \$US sur l'ensemble de la période, pour une valeur actuelle nette de 2\,122,7~millions \$US au taux d'actualisation de 10\,\%. Ce montant ne dépend pas du régime fiscal retenu~: il représente la « taille du gâteau » à partager entre l'État, l'investisseur et, le cas échéant, les créanciers.

Le tableau~\ref{tab:repartition} présente la répartition de ce flux entre l'État et l'investisseur sous deux régimes~: le régime « RDC amendée », qui demeure le régime actif par défaut de l'outil, et le régime « Sénégal » construit pour ce rapport. Sous le régime sénégalais, l'État perçoit 1\,589,7~millions \$US, soit 44,3\,\% du flux avant impôt, contre 2\,072,8~millions \$US (57,6\,\%) sous le régime congolais amendé --- un écart de 483~millions \$US sur la durée de vie du projet, qui revient symétriquement à l'investisseur.

\begin{table}[H]
\centering
\caption{Répartition du flux de trésorerie avant impôt entre l'État et l'investisseur, par régime (millions \$US réels, cumul 2020-2036)}
\label{tab:repartition}
\begin{tabular}{lccc}
\toprule
 & RDC -- amendée & Sénégal (SGO) & Écart \\
\midrule
Recettes de l'État          & 2\,072,8 & 1\,589,7 & $-$483,0 \\
Flux de l'investisseur après impôt & 1\,522,6 & 2\,005,6 & $+$483,0 \\
\midrule
Flux total avant impôt      & \multicolumn{2}{c}{3\,595,4} & -- \\
\TEMI                        & 57,6\,\% & 44,3\,\% & $-$13,3 points \\
\bottomrule
\end{tabular}
\end{table}

Cet écart de 13,3 points de \TEMI{} s'explique principalement par deux différences de paramétrage~: d'une part, le taux d'impôt sur les bénéfices, plus faible au Sénégal (25\,\% contre 35\,\%)~; d'autre part, l'exonération de droits de douane dont bénéficie SGO au titre de son statut d'entreprise franche d'exportation, qui n'a pas d'équivalent dans le régime congolais amendé (droits de douane de 5\,\%). À l'inverse, la redevance --- plus élevée au Sénégal (5\,\% contre 3,5\,\%, sur une assiette brute dans les deux cas) --- et la participation gratuite de l'État --- identique (10\,\%) --- jouent en sens inverse, mais de façon insuffisante pour compenser les deux premiers effets.

\section{Structure des recettes publiques par instrument fiscal}

Le tableau~\ref{tab:instruments} et la figure~\ref{fig:instruments} détaillent la composition des recettes de l'État pour les deux régimes, par instrument fiscal.

\begin{table}[H]
\centering
\caption{Recettes de l'État par instrument fiscal (millions \$US réels, cumul 2020-2036)}
\label{tab:instruments}
\small
\begin{tabular}{lrrrr}
\toprule
Instrument fiscal & \multicolumn{2}{c}{RDC -- amendée} & \multicolumn{2}{c}{Sénégal (SGO)} \\
 & Montant & Part & Montant & Part \\
\midrule
Bonus de signature                       & 0,0     & 0,0\,\%  & 0,0     & 0,0\,\%  \\
Droits de douane                         & 129,7   & 6,3\,\%  & 0,0     & 0,0\,\%  \\
Redevance                                & 239,9   & 11,6\,\% & 342,7   & 21,6\,\% \\
Impôt sur les bénéfices                  & 1\,121,1 & 54,1\,\% & 776,6  & 48,8\,\% \\
Impôt sur les profits excédentaires      & 224,9   & 10,9\,\% & 0,0     & 0,0\,\%  \\
Participation gratuite de l'État         & 188,0   & 9,1\,\%  & 247,6   & 15,6\,\% \\
Prélèvement à la source sur dividendes   & 169,2   & 8,2\,\%  & 222,8   & 14,0\,\% \\
\midrule
\textbf{Total} & \textbf{2\,072,8} & \textbf{100,0\,\%} & \textbf{1\,589,7} & \textbf{100,0\,\%} \\
\bottomrule
\end{tabular}

\vspace{0.2cm}
\footnotesize Les instruments « taxe provinciale », « redevance à taux variable sur les prix », « impôt sur la marge opérationnelle », « impôt sur la rente » et « impôt sur les bénéfices à taux variable » sont nuls dans les deux régimes et n'ont pas été reproduits dans le tableau.
\end{table}

\begin{figure}[H]
\centering
\begin{tikzpicture}
\begin{axis}[
  width=\textwidth, height=7.5cm,
  ybar, bar width=14pt,
  symbolic x coords={Droits de douane,Redevance,Impôt sur les bénéfices,Profits excédentaires,Participation État,Retenue dividendes},
  xtick=data, x tick label style={rotate=30,anchor=east,font=\scriptsize},
  ylabel={Millions \$US (réel, cumul 2020-2036)},
  ylabel style={font=\small}, tick label style={font=\small},
  enlarge x limits=0.12,
  legend style={at={(0.5,1.08)},anchor=south,legend columns=-1,font=\small},
  ymajorgrids, grid style={dashed,gray!30},
  axis lines*=left,
]
\addplot[fill=itieblue] coordinates {
 (Droits de douane,129.7)(Redevance,239.9)(Impôt sur les bénéfices,1121.1)
 (Profits excédentaires,224.9)(Participation État,188.0)(Retenue dividendes,169.2)
};
\addplot[fill=itiegold] coordinates {
 (Droits de douane,0)(Redevance,342.7)(Impôt sur les bénéfices,776.6)
 (Profits excédentaires,0)(Participation État,247.6)(Retenue dividendes,222.8)
};
\legend{RDC -- amendée, Sénégal (SGO)}
\end{axis}
\end{tikzpicture}
\caption{Recettes de l'État par instrument fiscal -- comparaison RDC amendée / Sénégal}
\label{fig:instruments}
\end{figure}

Trois enseignements se dégagent de cette comparaison. Premièrement, l'impôt sur les bénéfices demeure, dans les deux régimes, la principale source de recettes (49 à 54\,\% du total), ce qui confirme la sensibilité des recettes publiques à la profitabilité du projet --- donc au prix de l'or --- plutôt qu'à un instrument assis sur le volume. Deuxièmement, l'absence de droits de douane et d'impôt sur les profits excédentaires dans le régime sénégalais --- soit 17,2\,\% des recettes du régime congolais amendé --- n'est que partiellement compensée par le supplément de redevance (+102,8~millions \$US) et de participation de l'État (+59,6~millions \$US). Troisièmement, la retenue à la source sur les dividendes, dont le poids relatif passe de 8,2\,\% à 14,0\,\%, illustre le fait qu'une part accrue des recettes publiques sénégalaises est conditionnée à la distribution effective de dividendes par SGO --- recettes par construction plus tardives et plus incertaines que celles tirées de la redevance ou de l'impôt sur les bénéfices, qui sont prélevées dès que le projet est en exploitation et profitable, respectivement.

\section{Le \TEMI{} sénégalais dans la comparaison internationale}

La figure~\ref{fig:temi} et le tableau~\ref{tab:temi} positionnent le \TEMI{} obtenu pour le régime sénégalais --- 44,3\,\% --- parmi les dix régimes de comparaison déjà intégrés dans l'outil NRGI, tous appliqués au même profil de production que celui de Sabodala-Massawa. Cette mise en perspective permet d'apprécier, à profil opérationnel identique, l'effet du seul paramétrage fiscal sur le partage de la rente.

\begin{table}[H]
\centering
\caption{\TEMI{} par régime fiscal, appliqué au profil de production de Sabodala-Massawa}
\label{tab:temi}
\begin{tabular}{lc}
\toprule
Régime & \TEMI \\
\midrule
RDC -- amendée (2018)            & 57,6\,\% \\
Australie-Occidentale            & 54,7\,\% \\
Chili                             & 53,4\,\% \\
Ghana                             & 50,2\,\% \\
RDC -- taxe sur la rente         & 50,7\,\% \\
RDC -- redevance à taux variable & 44,8\,\% \\
\textbf{Sénégal (SGO)}           & \textbf{44,3\,\%} \\
RDC -- 2002                      & 44,5\,\% \\
Zambie -- 2016                   & 42,9\,\% \\
Tanzanie                         & 40,2\,\% \\
Pérou                             & 37,9\,\% \\
\bottomrule
\end{tabular}
\end{table}

\begin{figure}[H]
\centering
\begin{tikzpicture}
\begin{axis}[
  width=\textwidth, height=7.5cm,
  xbar, bar width=9pt,
  symbolic y coords={Pérou,Tanzanie,Zambie -- 2016,{Sénégal (SGO)},RDC -- 2002,{RDC -- red. variable},Ghana,{RDC -- taxe rente},Chili,{Australie-Occ.},{RDC -- amendée}},
  ytick=data, y tick label style={font=\scriptsize},
  xlabel={\TEMI{} (\%)}, xmin=0, xmax=65,
  xlabel style={font=\small}, tick label style={font=\small},
  enlarge y limits=0.07,
  xmajorgrids, grid style={dashed,gray!30},
  axis lines*=left,
  nodes near coords={\pgfmathprintnumber[fixed,fixed zerofill,precision=1]{\pgfplotspointmeta}},
  every node near coord/.append style={font=\scriptsize, xshift=6pt},
]
\addplot[fill=itieblue] coordinates {
 (37.9,Pérou)
 (40.2,Tanzanie)
 (42.9,Zambie -- 2016)
 (44.3,{Sénégal (SGO)})
 (44.5,RDC -- 2002)
 (44.8,{RDC -- red. variable})
 (50.2,Ghana)
 (50.7,{RDC -- taxe rente})
 (53.4,Chili)
 (54.7,{Australie-Occ.})
 (57.6,{RDC -- amendée})
};
\end{axis}
\end{tikzpicture}
\caption{Comparaison internationale du \TEMI{}, profil de production Sabodala-Massawa (le Sénégal est représenté en quatrième position à partir du bas)}
\label{fig:temi}
\end{figure}

Le \TEMI{} sénégalais se situe ainsi dans la partie inférieure de la distribution des onze régimes comparés, à un niveau quasiment identique à celui du code minier congolais de 2002 (avant l'amendement de 2018) et au régime congolais à redevance variable, mais nettement inférieur à celui du code minier congolais amendé, du régime australien et du régime chilien --- pourtant souvent cités comme exemples de régimes compétitifs pour les pays développés. Deux régimes seulement --- la Tanzanie et le Pérou --- présentent un \TEMI{} inférieur à celui du Sénégal pour ce profil de projet. Cette position relativement basse dans la comparaison internationale doit cependant être interprétée à la lumière de la remarque formulée au chapitre~\ref{chap:methodo}~: le régime sénégalais modélisé ne couvre pas les flux additionnels liés à la transaction Massawa (paiement initial et \emph{gold stream}), qui --- bien que de nature contractuelle et non fiscale --- représentent, sur la durée du PFS, un ordre de grandeur de 155~millions \$US (15~millions \$US de paiement initial et environ 140~millions \$US de valeur actualisée du \emph{gold stream}), soit près de 10\,\% des recettes fiscales calculées par le modèle pour le régime sénégalais. Leur prise en compte rapprocherait le \TEMI{} effectif de SGO du niveau observé pour le régime congolais à taxe sur la rente ou pour le Ghana.

\section{Limites de l'approche pour un projet d'extension de mine}
\label{sec:limites-tri}

Les indicateurs de rentabilité du point de vue de l'investisseur calculés par l'outil --- \TRI{} avant et après impôt, période de récupération actualisée et non actualisée --- renvoient tous une valeur nulle, quel que soit le régime fiscal retenu (tableau~\ref{tab:tri}).

\begin{table}[H]
\centering
\caption{Indicateurs de rentabilité calculés par le modèle}
\label{tab:tri}
\begin{tabular}{lc}
\toprule
Indicateur & Valeur \\
\midrule
\VAN{} avant impôt (actualisée à 10\,\%) & 2\,122,7 millions \$US \\
\TRI{} avant impôt & non défini (0 affiché) \\
\TRI{} après impôt & non défini (0 affiché) \\
Période de récupération (non actualisée) & non définie (0 affiché) \\
Période de récupération (actualisée) & non définie (0 affiché) \\
\bottomrule
\end{tabular}
\end{table}

Cette absence de résultat n'est pas une erreur de calcul, mais la conséquence directe de la nature du profil de Sabodala-Massawa~: le flux de trésorerie avant impôt de la première année du modèle (2020) est déjà positif (143,0~millions \$US), et il le reste pour chacune des dix-sept années du projet. Or, par construction, le \TRI{} d'une série de flux de trésorerie n'existe que si cette série change au moins une fois de signe~: la fonction financière utilisée par le modèle ne trouve donc aucune racine et renvoie une erreur, que la formule convertit par défaut en zéro. La période de récupération de l'investissement --- définie comme le nombre d'années nécessaire pour que le cumul des flux de trésorerie actualisés devienne positif --- est, pour la même raison, immédiatement atteinte (dès l'année~0) et n'est donc pas non plus interprétable comme un indicateur de risque.

Cette situation est caractéristique des outils FARI/NRGI, conçus à l'origine pour des projets \emph{greenfield}, dont la phase de construction se traduit par plusieurs années de flux négatifs avant le démarrage de la production. Sabodala-Massawa, en tant qu'\textbf{extension} d'une mine déjà en exploitation --- où les dépenses d'investissement de la phase~1 de Massawa (160~millions \$US en 2022, soit la plus importante annuité du profil) sont financées par les flux de trésorerie dégagés par les opérations existantes --- ne présente, à l'échelle du modèle, aucune année de flux net négatif.

Cette limite n'invalide pas pour autant les résultats présentés aux sections précédentes~: la \VAN{} et le \TEMI, qui ne reposent pas sur une condition de changement de signe, restent parfaitement définis et interprétables. Elle invite néanmoins à la prudence quant à l'utilisation de ce type d'outil pour les décisions d'investissement relatives à des projets d'extension~: un investisseur ou une administration souhaitant évaluer la rentabilité \emph{marginale} de la décision d'investissement (ici, la mise en valeur du gisement Massawa) devrait isoler, dans le flux de trésorerie, la seule part additionnelle liée à cette décision --- ce qui sortirait du cadre du présent exercice, fondé sur le profil consolidé du complexe Sabodala-Massawa tel que présenté dans le PFS.
